\documentclass{amsart}
% For dealing with references we use the comment environment
\usepackage{verbatim}
\newenvironment{reference}{\comment}{\endcomment}
%\newenvironment{reference}{}{}
\newenvironment{slogan}{\comment}{\endcomment}
\newenvironment{history}{\comment}{\endcomment}

% For commutative diagrams we use Xy-pic
\usepackage[all]{xy}

% We use 2cell for 2-commutative diagrams.
\xyoption{2cell}
\UseAllTwocells

% We use multicol for the list of chapters between chapters
\usepackage{multicol}

% This is generally recommended for better output
\usepackage{lmodern}
\usepackage[T1]{fontenc}

% For cross-file-references

% Package for hypertext links:
\usepackage{hyperref}

% For any local file, say "hello.tex" you want to link to please

% Theorem environments.
%
\theoremstyle{plain}
\newtheorem{theorem}[subsection]{Theorem}
\newtheorem{proposition}[subsection]{Proposition}
\newtheorem{lemma}[subsection]{Lemma}

\theoremstyle{definition}
\newtheorem{definition}[subsection]{Definition}
\newtheorem{example}[subsection]{Example}
\newtheorem{exercise}[subsection]{Exercise}
\newtheorem{situation}[subsection]{Situation}

\theoremstyle{remark}
\newtheorem{remark}[subsection]{Remark}
\newtheorem{remarks}[subsection]{Remarks}

\numberwithin{equation}{subsection}

% Macros
%
\def\lim{\mathop{\mathrm{lim}}\nolimits}
\def\colim{\mathop{\mathrm{colim}}\nolimits}
\def\Spec{\mathop{\mathrm{Spec}}}
\def\Hom{\mathop{\mathrm{Hom}}\nolimits}
\def\Ext{\mathop{\mathrm{Ext}}\nolimits}
\def\SheafHom{\mathop{\mathcal{H}\!\mathit{om}}\nolimits}
\def\SheafExt{\mathop{\mathcal{E}\!\mathit{xt}}\nolimits}
\def\Sch{\mathit{Sch}}
\def\Mor{\mathop{\mathrm{Mor}}\nolimits}
\def\Ob{\mathop{\mathrm{Ob}}\nolimits}
\def\Sh{\mathop{\mathit{Sh}}\nolimits}
\def\NL{\mathop{N\!L}\nolimits}
\def\CH{\mathop{\mathrm{CH}}\nolimits}
\def\proetale{{pro\text{-}\acute{e}tale}}
\def\etale{{\acute{e}tale}}
\def\QCoh{\mathit{QCoh}}
\def\Ker{\mathop{\mathrm{Ker}}}
\def\Im{\mathop{\mathrm{Im}}}
\def\Coker{\mathop{\mathrm{Coker}}}
\def\Coim{\mathop{\mathrm{Coim}}}

% Boxtimes
%
\DeclareMathSymbol{\boxtimes}{\mathbin}{AMSa}{"02}

%
% Macros for moduli stacks/spaces
%
\def\QCohstack{\mathcal{QC}\!\mathit{oh}}
\def\Cohstack{\mathcal{C}\!\mathit{oh}}
\def\Spacesstack{\mathcal{S}\!\mathit{paces}}
\def\Quotfunctor{\mathrm{Quot}}
\def\Hilbfunctor{\mathrm{Hilb}}
\def\Curvesstack{\mathcal{C}\!\mathit{urves}}
\def\Polarizedstack{\mathcal{P}\!\mathit{olarized}}
\def\Complexesstack{\mathcal{C}\!\mathit{omplexes}}
% \Pic is the operator that assigns to X its picard group, usage \Pic(X)
% \Picardstack_{X/B} denotes the Picard stack of X over B
% \Picardfunctor_{X/B} denotes the Picard functor of X over B
\def\Pic{\mathop{\mathrm{Pic}}\nolimits}
\def\Picardstack{\mathcal{P}\!\mathit{ic}}
\def\Picardfunctor{\mathrm{Pic}}
\def\Deformationcategory{\mathcal{D}\!\mathit{ef}}

\setlength{\emergencystretch}{3em}
\begin{document}
\title[Stacks Project, Tag 0G9K]{405 --- Power annihilation of Ext and local projectivity of a two-term derived object}
\maketitle
\noindent Copyright (C) 2005--2025 Johan de Jong. Stacks Project authors. Source: \href{https://stacks.math.columbia.edu/tag/0G9K}{Tag 0G9K}.

\noindent Editorial difficulty note (not author text). Difficulty H2: The author uses a cohomological triangle and injective targets to separate the negative cohomology obstruction from near-projectivity of degree-zero cohomology. A localized universal extension class and uniform radical-power estimates recover the converse from projectivity off the ideal locus..

\noindent Editorial provenance note (not author text). History: excerpt from revision \texttt{a0444\allowbreak 6e57e\allowbreak c1fbc\allowbreak 252a8\allowbreak 71afc\allowbreak ec775\allowbreak 2fb28\allowbreak 07b14}, more-algebra.tex, lines 22629--22709. Reference presentation was adapted; original-author context and core support blocks were retained or added. The mathematical wording is unchanged. Licensed under GNU FDL 1.2 or later; no invariant sections or cover texts. See ../licenses/COPYING. Context blocks reproduce author definitions and supporting results; they are not additional counted problems.

\begin{lemma}
\label{lemma-ext-1-annihilated-definite}
Let $R$ be a ring and let $I \subset R$ be an ideal.
Let $K \in D(R)$. Assume $H^i(K) = 0$ for $i \not \in \{-1, 0\}$.
The following are equivalent
\begin{enumerate}
\item $\Ext^1_R(K, N)$ is annihilated by $I$ for all $R$-modules $N$,
\item $K$ can be represented by a complex $K^{-1} \to K^0$
with $K^0$ free such that for any $a \in I$ the map
$a : K^{-1} \to K^{-1}$ factors through $d_K^{-1} : K^{-1} \to K^0$,
\item whenever $K$ is represented by a two term complex
$K^{-1} \to K^0$ with $K^0$ projective, then for any $a \in I$ the map
$a : K^{-1} \to K^{-1}$ factors through $d_K^{-1} : K^{-1} \to K^0$.
\end{enumerate}
If $R$ is Noetherian and $H^i(K)$ is a finite $R$-module for $i = -1, 0$,
then these are also equivalent to
\begin{enumerate}
\item[(4)] $\Ext^1_R(K, N)$ is annihilated by $I$ for every finite
$R$-module $N$,
\item[(5)] $K$ can be represented by a complex $K^{-1} \to K^0$
with $K^0$ finite free and $K^{-1}$ finite such that for any $a \in I$ the map
$a : K^{-1} \to K^{-1}$ factors through $d_K^{-1} : K^{-1} \to K^0$.
\end{enumerate}
\end{lemma}

\begin{definition}
\label{definition-near-projective}
Let $R$ be a ring. Let $I \subset R$ be an ideal. Let $M$ be an $R$-module.
We say $M$ is {\it $I$-projective}\footnote{This is nonstandard notation.}
if the equivalent conditions of Lemma \ref{lemma-near-projective} hold.
\end{definition}

\begin{lemma}
\label{lemma-near-projective}
Let $R$ be a ring. Let $I \subset R$ be an ideal. Let $M$ be an $R$-module.
The following conditions are equivalent
\begin{enumerate}
\item for every $a \in I$ the map $a : M \to M$ factors through a projective
$R$-module,
\item for every $a \in I$ the map $a : M \to M$ factors through a free
$R$-module, and
\item $\Ext^1_R(M, N)$ is annihilated by $I$ for every $R$-module $N$.
\end{enumerate}
\end{lemma}

\begin{lemma}
\label{lemma-ext-1-annihilated}
Let $R$ be ring and let $I \subset R$ be an ideal.
Let $K \in D(R)$ with $H^i(K) = 0$ for $i \not \in \{-1, 0\}$.
The following are equivalent
\begin{enumerate}
\item there exists a $c \geq 0$ such that the equivalent
conditions (1), (2), (3) of Lemma \ref{lemma-ext-1-annihilated-definite}
hold for $K$ and the ideal $I^c$,
\item there exists a $c \geq 0$ such that (a) $I^c$ annihilates
$H^{-1}(K)$ and (b) $H^0(K)$ is an $I^c$-projective module (see
Section \href{https://stacks.math.columbia.edu/tag/0G8Z}{section near projective (Tag 0G8Z)}).
\end{enumerate}
If $R$ is Noetherian and $H^i(K)$ is a finite $R$-module
for $i = -1, 0$, then these are also equivalent to
\begin{enumerate}
\item[(3)] there exists a $c \geq 0$ such that the equivalent
conditions (4), (5) of Lemma \ref{lemma-ext-1-annihilated-definite}
hold for $K$ and the ideal $I^c$,
\item[(4)] $H^{-1}(K)$ is $I$-power torsion and there exist
$f_1, \ldots, f_s \in R$ with $V(f_1, \ldots, f_s) \subset V(I)$
such that the localizations $H^0(K)_{f_i}$ are projective
$R_{f_i}$-modules,
\item[(5)] $H^{-1}(K)$ is $I$-power torsion and there exist
$f_1, \ldots, f_s \in I$ with $V(f_1, \ldots, f_s) = V(I)$
such that the localizations $H^0(K)_{f_i}$ are projective
$R_{f_i}$-modules, and
\item[(6)] $H^{-1}(K)$ is $I$-power torsion and for any
$f_1, \ldots, f_s \in I$ with $V(f_1, \ldots, f_s) = V(I)$
the localizations $H^0(K)_{f_i}$ are projective $R_{f_i}$-modules.
\end{enumerate}
\end{lemma}

\begin{proof}
The distinguished triangle $H^{-1}(K)[1] \to K \to H^0(K)[0] \to H^{-1}(K)[2]$
determines an exact sequence
$$
0 \to \Ext^1_R(H^0(K), N) \to \Ext^1_R(K, N) \to \Hom_R(H^{-1}(K), N) \to
\Ext^2_R(H^0(K), N)
$$
Thus (2) implies that $I^{2c}$ annihilates $\Ext^1_R(K, N)$ for every
$R$-module $N$. Assuming (1) we immediately see that $H^0(K)$ is
$I^c$-projective. On the other hand, we may choose an injective map
$H^{-1}(K) \to N$ for some injective $R$-module $N$. Then this map
is the image of an element of $\Ext^1_R(K, N)$ by the vanishing
of the $\Ext^2$ in the sequence and we conclude $H^{-1}(K)$
is annihilated by $I^c$.

\medskip\noindent
Assume $R$ is Noetherian and $H^i(K)$ is a finite $R$-module
for $i = -1, 0$. By Lemma \ref{lemma-ext-1-annihilated-definite}
we see that (3) is equivalent to (1) and (2). Also, if (3)
holds then for $f \in I$ the multiplication by $f$ on
$H^0(K)$ factors through a projective module, which implies
that $H^0(K)_f$ is a summand of a projective $R_f$-module and hence
itself a projective $R_f$-module. Thus the equivalent conditions
(1), (2), and (3) imply (6).
Of course (6) implies (5) and (5) trivially implies (4).

\medskip\noindent
Assume (4). Since $H^{-1}(K)$ is a finite $R$-module and $I$-power torsion
we see that $I^{c_1}$ annihilates $H^{-1}(K)$ for some $c_1 \geq 0$.
Choose a short exact sequence
$$
0 \to M \to R^{\oplus r} \to H^0(K) \to 0
$$
which determines an element $\xi \in \Ext^1_R(H^0(K), M)$.
For any $f \in I$ we have $\Ext^1_R(H^0(K), M)_f = \Ext^1_{R_f}(H^0(K)_f, M_f)$
by Lemma \href{https://stacks.math.columbia.edu/tag/087R}{lemma pseudo coherence and base change ext (Tag 087R)}.
Hence if $H^0(K)_f$ is projective, then a power of $f$ annihilates $\xi$.
We conclude that $\xi$ is annihilated by $(f_1, \ldots, f_s)^{c_2}$
for some $c_2 \geq 0$. Since $V(f_1, \ldots, f_s) \subset V(I)$ we have
$\sqrt{I} \subset (f_1, \ldots, f_s)$
(Algebra, Lemma \href{https://stacks.math.columbia.edu/tag/00E0}{lemma Zariski topology (Tag 00E0)}).
Since $R$ is Noetherian we find $I^{c_3} \subset (f_1, \ldots, f_s)$
for some $c_3 \geq 0$ (Algebra, Lemma \href{https://stacks.math.columbia.edu/tag/00IM}{lemma Noetherian power (Tag 00IM)}).
Hence $I^{c2c3}$ annihilates $\xi$.
This in turn says that $H^0(K)$ is $I^{c_2c_3}$-projective (as multiplication
by $a \in I$ which annihilate $\xi$ factor through $R^{\oplus r}$).
Hence taking $c = \max(c_1, c_2c_3)$ we see that (2) holds.
\end{proof}
\end{document}
